% ===== PRÉAMBULE CANONIQUE — à coller VERBATIM en tête de CHAQUE .tex =====
\documentclass[11pt,a4paper]{article}
\usepackage[utf8]{inputenc}\usepackage[T1]{fontenc}\usepackage{lmodern}
\usepackage[french]{babel}
\usepackage{amsmath,amssymb,mathtools}\usepackage{icomma}
\usepackage{siunitx}\sisetup{locale=FR}  % \qty{}{}, \si{}, \ang{}
\usepackage{xcolor}\usepackage{colortbl}
\usepackage{tikz}\usepackage{pgfplots}\pgfplotsset{compat=1.18}
\usepackage[siunitx,europeanresistors]{circuitikz} % circuits (élec.)
\usepackage[most]{tcolorbox}\usepackage{enumitem}\usepackage{array,booktabs}
\usepackage[a4paper,margin=2.2cm]{geometry}
\usetikzlibrary{arrows.meta,calc,angles,quotes,positioning,patterns,%
  backgrounds,shapes.geometric,decorations.pathmorphing,decorations.markings,intersections}
\definecolor{primary}{RGB}{222,49,99}\definecolor{primarydark}{RGB}{150,22,55}
\definecolor{defcol}{RGB}{0,114,178}\definecolor{methcol}{RGB}{52,92,148}
\definecolor{excol}{RGB}{0,135,90}\definecolor{attncol}{RGB}{200,40,55}
\tcbset{boxbase/.style={enhanced,breakable,boxrule=0.9pt,arc=2.5pt,
  left=3mm,right=3mm,top=2.3mm,bottom=2.3mm,fonttitle=\bfseries,coltitle=white}}
% --- boîtes TUTORIEL (noms EXACTS, ne pas renommer) ---
\newtcolorbox{cle}[1][]{boxbase,colback=defcol!6,colframe=defcol,
  title={\textsc{À retenir}\ifx\relax#1\relax\relax\else\textmd{ : \itshape #1}\fi}}
\newtcolorbox{methode}[1][]{boxbase,colback=methcol!7,colframe=methcol,sharp corners,
  title={\textsc{Méthode}\ifx\relax#1\relax\relax\else\textmd{ --- \itshape #1}\fi}}
\newtcolorbox{exemplebox}[1][]{boxbase,colback=excol!5,colframe=excol,
  title={\textsc{Exemple}\ifx\relax#1\relax\relax\else\textmd{ : \itshape #1}\fi}}
\newtcolorbox{piege}[1][]{boxbase,colback=attncol!6,colframe=attncol,
  title={\textsc{Piège}\ifx\relax#1\relax\relax\else\textmd{ : \itshape #1}\fi}}
\newtcolorbox{remarque}[1][]{boxbase,colback=black!4,colframe=black!45,
  title={\textsc{Remarque}\ifx\relax#1\relax\relax\else\textmd{ : \itshape #1}\fi}}
% --- boîtes EXERCICE / CORRIGÉ (noms EXACTS, auto-comptées) ---
\newtcolorbox[auto counter]{exo}[1][]{boxbase,colback=primary!4,colframe=primary,
  title={\textsc{Exercice}~\thetcbcounter\ifx\relax#1\relax\relax\else\textmd{ --- \itshape #1}\fi}}
\newtcolorbox[auto counter]{corrige}[1][]{boxbase,colback=excol!4,colframe=excol,
  title={\textsc{Corrigé}~\thetcbcounter\ifx\relax#1\relax\relax\else\textmd{ --- \itshape #1}\fi}}
\newcommand{\vv}[1]{\vec{#1}}\newcommand{\ds}{\displaystyle}
\newcommand{\R}{\mathbb{R}}\newcommand{\N}{\mathbb{N}}\newcommand{\Z}{\mathbb{Z}}
% (un \usepackage{fancyhdr}+en-tête est possible ; il est ignoré côté web)
% =========================================================================
\begin{document}
\begin{corrige}[calculer la force de Laplace]
Le conducteur est perpendiculaire à $\vv{B}$ ($\theta=\ang{90}$, $\sin\theta=1$), donc $F=BIL$ :
\[ F=B\,I\,L=0{,}4\times3\times0{,}15=\qty{0.18}{\newton}. \]
(Bien convertir $L=\qty{15}{\centi\metre}=\qty{0.15}{\metre}$.)
\end{corrige}

\begin{corrige}[le conducteur parallèle au champ]
Le conducteur est \textbf{parallèle} à $\vv{B}$, donc $\theta=\ang{0}$ et $\sin\ang{0}=0$ :
\[ F=B\,I\,L\,\sin\ang{0}=\qty{0}{\newton}. \]
La force de Laplace est \textbf{nulle} : un conducteur parallèle au champ ne subit aucune force,
\emph{quelle que soit} l'intensité du courant. C'est le piège classique du concours.
\end{corrige}

\begin{corrige}[trouver le sens de la force]
Règle des trois doigts (main droite) : \textbf{pouce} vers la droite (sens de $I$), \textbf{index}
vers le haut (sens de $\vv{B}$) ; le \textbf{majeur}, perpendiculaire au plan, pointe alors
\textbf{vers toi}. La force \textbf{sort de la feuille}, $\odot$.
\end{corrige}

\begin{corrige}[deux fils, même sens de courant]
Deux courants de \textbf{même sens} $\Rightarrow$ les fils \textbf{s'attirent}.
(On peut le retrouver : le fil 1 crée en 2 un champ $\vv{B}_1$ ; la force de Laplace sur le fil 2,
par la règle des trois doigts, est dirigée vers le fil 1. Réciproquement pour le fil 1.)
\end{corrige}

\begin{corrige}[l'effet de l'angle]
La force varie comme $\sin\theta$ : $F=F_0\,\dfrac{\sin\theta}{\sin\ang{90}}=F_0\sin\theta$.
\begin{enumerate}[label=\alph*)]
  \item $F=F_0\sin\ang{30}=0{,}2\times0{,}5=\qty{0.1}{\newton}$ (la moitié).
  \item $F=F_0\sin\ang{0}=\qty{0}{\newton}$ : parallèle au champ, force nulle.
\end{enumerate}
\end{corrige}

\end{document}
